\section{Geometric grids and the Fourier upper bound}\label{sec:last}

The character table of a finite abelian group factors into cyclic
character tables. The following result allows arbitrary geometric
progressions on the axes, allows interpolation in place of evaluation,
and admits triangular Toeplitz factors in the same product. Only the
prescribed factors must be defined over the base field. Auxiliary
Fourier transforms will take place in an extension whose coordinates
hold independent fibers of the original input.

\begin{theorem}[Geometric grids and triangular Toeplitz factors]\label{up:grid}\label{up:fourier}
Let \(K\) be any field, \(a_1,\ldots,a_r\ge2\), and \(N=\prod_i a_i\).
For each \(i\), let \(T_i\) be a matrix of side \(a_i\) over \(K\).
It may be the geometric evaluation matrix
\(V_{a_i}(q_i)=(q_i^{jk})_{0\le j,k<a_i}\) at distinct nodes with
\(q_i\in K^\times\), its inverse, or an invertible upper or lower
triangular Toeplitz matrix. For every \(\alpha>1/2\), the operator
\(\bigotimes_i T_i\) has a binary linear circuit over \(K\) of size
\(O_\alpha(N(1+\log N)^\alpha)\), with an implied constant uniform
in the field, the dimensions and the factors.
\end{theorem}

\begin{proof}
Fix \(1/2<\beta<\min(\alpha,1)\), put \(t=\lceil\log_2 N\rceil\) and
take \(c=1/24\). An axis of length \(a>N^c\) can be handled individually.
Corollary~\ref{cor:toe} applies a triangular Toeplitz factor on all
\(N/a\) fibers for \(O_\beta(N(1+t)^\beta)\) gates. By the Newton
factorization (Lemma~\ref{up:newton-lemma}), geometric evaluation and
its inverse satisfy the same bound. There are at most \(1/c\) such axes.

Group the smaller axes by multiplying their lengths until the product
first reaches \(N^c\), then start a new group. The last group may
fall short. There are \(O(1/c)\) groups and each has volume
\(B\le N^{2c}\). Fix one group. The \(h=N/B\) choices of complementary
coordinates supply \(h\) independent arrays, each of size \(B\).
Our two budgets for the group are a word of length \(O(1+\log t)\)
on \(O(B)\) registers, and an extension of degree at most \(h\).
The first will bound the number of passes through these arrays.
The second will let us fill the extension coordinates with data.

\paragraph{Choosing auxiliary lengths.}
A Fourier word for one factor of size \(a\le N^c\) needs a working
length \(b_a\) between \(2a\) and \(6a\) (Lemma~\ref{up:toeplitz}).
Directly padding to a power of the radix \(\ell\) could introduce
\(O(\log a)\) successive stages. We choose \(b_a\) with pairwise
coprime factors of size \(O(t^2)\). Their cyclic transforms can then
be synchronized and only the short factors need radix padding.

Use \(\ell=2\) except in characteristic two, where \(\ell=3\).
List the primes in \([t,t^2]\), omitting \(\ell\) and the
characteristic. By the prime number theorem their product exceeds
\(2N\) for large \(t\), even after the omissions. Let \(P\) be the
product of the longest prefix with \(P<2a\), taking \(P=1\) for the
empty prefix. Set \(b_a=P\ell^j\) for the least \(j\ge0\) with
\(P\ell^j\ge2a\). Then \(2a\le b_a<6a\). The next available prime
would have sufficed in place of \(\ell^j\), so \(\ell^j\le3t^2\).
The primes in \(P\) and this radix power give pairwise coprime
factors \(d\le3t^2\). For each \(d\ge2\), take the least power \(s_d\)
of \(\ell\) with \(s_d\ge2d\), so \(2d\le s_d<6d\).

Let \(M\) be the least common multiple of the lengths \(b_a,s_d\)
in this group and set \(E=K(\zeta_M)\) for a primitive \(M\)th root.
Every working length is prime to the characteristic and divides \(M\).
All the \(s_d\) are powers of one radix and their maximum is \(O(t^2)\).
If the group has \(k\) axes, then \(2^k\le B\), giving
\(\prod_a b_a\le6^k B\le B^4\), with one factor per axis. Hence
\begin{equation}\label{up:capacity}
               e=[E:K]\le M=O(t^2B^4)\le N/B=h
\end{equation}
for all sufficiently large \(N\), since \(B\le N^{1/12}\).

\paragraph{Constructing the word.}
Begin with the radix Fourier word of length \(O(1+\log t)\) on
\(s_d=O(t^2)\) coordinates from Section~\ref{sec:lay}.
Lemma~\ref{up:toeplitz} with \(a=d\), \(b=s_d\) gives a word of the same
order of length for \(F_d\), with \(O(d)\) registers.
Good's prime-factor factorization \cite{Good} expresses \(F_{b_a}\)
as the tensor product of these short transforms at appropriate primitive
roots, up to the Chinese remainder permutations on the input and output
indices. Lemma~\ref{up:fibers} synchronizes them within the same length
bound using \(O(b_a)\) registers.

Apply Lemma~\ref{up:toeplitz} a second time, now using \(F_{b_a}\) to
implement the prescribed factor \(T_i\) with \(a_i=a\). Synchronize these
axis words by Lemma~\ref{up:fibers}. The result is the desired word
of length \(O(1+\log t)\) on \(O(B)\) registers, acting as the tensor
product of the prescribed factors and restoring its other registers.
The original sizes \(a_i\) may be divisible by the characteristic.
Only the auxiliary Fourier lengths had to avoid it.

\paragraph{Returning to the base field.}\label{up:packing}
Let \(T\) be the transform of this group of axes. Pack \(e\) independent
input arrays \(x_0,\ldots,x_{e-1}\) entrywise in the power basis of
\(\zeta_M\). The intermediate word uses \(E\)-scalars, but \(T\) itself
has entries in \(K\). Thus each base coordinate returns its own answer,
\[
       T\left(\sum_{\nu=0}^{e-1}\zeta_M^\nu x_\nu\right)
        =\sum_{\nu=0}^{e-1}\zeta_M^\nu T x_\nu .
\]
By \eqref{up:capacity}, \(\lceil h/e\rceil\le2h/e\) batches suffice,
using zeros for unused coordinates of the last batch.

One monomial layer permutes or scales \(O(B)\) extension entries per
batch. Multiplication by a prescribed \(E\)-scalar costs
\(O_\beta(e(1+\log e)^\beta)\) binary linear gates over \(K\). Indeed,
multiply in the power basis and reduce modulo the fixed monic minimal
polynomial of \(\zeta_M\). The reversed quotient is a truncated product
of the reversed high coefficients with the reciprocal of the reversed
divisor. Subtracting the product of quotient and divisor gives the
remainder. Each convolution has one fixed operand, so
Corollary~\ref{cor:toe} supplies the stated bound. Over all batches
the layer costs \(O_\beta(N(1+t)^\beta)\), since \(e\le h\le N\).

A shear layer has blocks of total size \(O(B)\). On a block of size
\(2^j\) the entries of \(U^{\otimes j}\) lie in the prime field.
The block therefore acts independently on the \(e\) base coordinates.
By Lemma~\ref{lem:shear}, each costs \(O_\beta(2^j(1+j)^\beta)\)
gates. Since \(j=O(1+t)\), the entire shear layer also costs
\(O_\beta(N(1+t)^\beta)\) over all batches. There are \(O(1+\log t)\)
layers per group and boundedly many groups. The factor \(1+\log t\)
is \(O_{\alpha,\beta}((1+t)^{\alpha-\beta})\), which yields the claimed
bound. Direct circuits cover the finitely many smaller \(N\).
\end{proof}

\begin{proof}[Completion of the proof of Theorem~\ref{thm:main}]
Write a group as \(G=\prod_i C_{a_i}\), and choose a primitive
\(a_i\)th root \(q_i\) in the given splitting field. Its character
transform is \(\bigotimes_i V_{a_i}(q_i)\), up to the order of the
characters and group elements. Its inverse is the tensor product of the inverse factors. Theorem~\ref{up:grid} supplies a circuit
of size \(O_\alpha(N(1+\log N)^\alpha)\) for either choice. Applying
these circuits separately to the \(J\) independent blocks gives the
upper bound with one constant for all \(J\), all groups of the given
order and all common splitting fields. Together with the lower bound
this determines the exponent.
\end{proof}

\subsection{The logarithmic exponent of formal convolution}

Fixing one operand of a polynomial product makes the output linear in
the other operand. A suitable choice puts a cyclic Fourier transform
among the output coefficients. Extracting the linear part at every gate
then transfers the Fourier lower bound to division-free circuits for
multiplication.

\begin{corollary}[Independent formal convolutions]\label{cor:conv-exponent}
For every field \(K\), \(n\ge2\) and \(J\ge1\), let \(S_K(J,n)\) be
the minimum size of a binary division-free arithmetic circuit computing
\(J\) independent products of polynomials of degree less than \(n\).
All coefficients of all products are required as formal outputs. Then
there is an absolute \(c>0\) and, for each \(\alpha>1/2\), a constant
\(C_\alpha\) such that
\[
       cJn\sqrt{\log n}\ \le\ S_K(J,n)\
                  \le\ C_\alpha Jn(1+\log n)^\alpha .
\]
In particular the logarithmic exponent for formal convolution is
\(1/2\), uniformly in the field and the number of products.
\end{corollary}

\begin{proof}
Theorem~\ref{thm:conv} supplies the upper bound by computing each
product separately. For the lower bound, extend a size-\(s\) circuit
from \(K\) to \(\overline K\). For \(n\ge8\), choose
\(m=\lfloor n/2\rfloor\) or \(m=\lfloor n/2\rfloor-1\) invertible in
this field, and let \(q\) be a primitive \(m\)th root. Both lengths
are comparable to \(n\), and at most one is divisible by the
characteristic.

Bluestein's chirp identity \cite{Bluestein} uses the integer
\(b(k)=k(k-1)/2\), which satisfies \(b(i)+b(-j)-b(i-j)=ij\).
In each product set \(y_j=q^{b(-j)}x_j\) for \(0\le j<m\), using an
independent vector \(x\) for each block. Fix the other operand in every
block to the same polynomial \(H\), with coefficients
\(H_{m-1+d}=q^{-b(d)}\) for \(|d|<m\). Its \(2m-1\) coefficients fit
in length \(n\). Padding with zeros as needed gives, in each block,
\[
  q^{b(i)}[X^{m-1+i}]\left(H(X)\sum_{j=0}^{m-1} y_j X^j\right)
       =\sum_{j=0}^{m-1} q^{ij}x_j,\qquad 0\le i<m .
\]

Take the homogeneous linear part at each gate in \emph{all} remaining
variables together. Sums and differences combine the linear parts of
their arguments. At a product \(uv\) the linear part is
\(u(0)v_1+v(0)u_1\), using one binary linear gate for the linear
parts \(u_1,v_1\), even if both depend on many blocks. Input and output
scalings in the chirp identity can be absorbed into coefficients and
readouts. Thus a linear circuit with at most \(s\) gates computes
\(J\) independent Fourier transforms on \(C_m\). Theorem~\ref{thm:main}
gives \(s=\Omega(Jm\sqrt{\log m})=\Omega(Jn\sqrt{\log n})\).
For the finitely many smaller \(n\), the \(J(2n-1)\) output
coefficients are distinct nonzero quadratic forms. Each needs a gate
output distinct up to scalars from those for the others, which gives
the required linear bound in \(Jn\). The constants are uniform in
\(J\) and \(K\).
\end{proof}
